\begin{register}{H}{GPIO\_OUT\_REG}{0x0004}\label{regdesc:GPIOOUTREG}
\regfield{GPIO\_OUT\_DATA\_ORIG}{32}{0}{{0x000000}}%
\reglabel{Reset}\regnewline%
\begin{regdesc}\begin{reglist}
\label{fielddesc:GPIOOUTDATAORIG}\item [GPIO\_OUT\_DATA\_ORIG] 配置简单 GPIO 输出模式下 GPIO0 \verb+~+ GPIO31 的输出值。\\
0：低电平\\
1：高电平 \\
 bit0 \verb+~+ bit31 的值分别对应 GPIO0 \verb+~+ 31 的输出值。\\
(R/W/SC/WTC)
\end{reglist}\end{regdesc}
\end{register}


\begin{register}{H}{GPIO\_OUT\_W1TS\_REG}{0x0008}\label{regdesc:GPIOOUTW1TSREG}
\regfield{GPIO\_OUT\_W1TS}{32}{0}{{0x000000}}%
\reglabel{Reset}\regnewline%
\begin{regdesc}\begin{reglist}
\label{fielddesc:GPIOOUTW1TS}\item [GPIO\_OUT\_W1TS] 配置是否置位 GPIO0 \verb+~+ GPIO31 的输出寄存器 \hyperref[regdesc:GPIOOUTREG]{GPIO\_OUT\_REG}。\\
0:不置位\\
1： \hyperref[regdesc:GPIOOUTREG]{GPIO\_OUT\_REG} 中相应位也将置 1 \\
bit0 \verb+~+ bit31 分别对应 GPIO0 \verb+~+ GPIO31。
推荐使用此寄存器来置位  \hyperref[regdesc:GPIOOUTREG]{GPIO\_OUT\_REG}。\\
(WT)
\end{reglist}\end{regdesc}
\end{register}


\begin{register}{H}{GPIO\_OUT\_W1TC\_REG}{0x000C}\label{regdesc:GPIOOUTW1TCREG}
\regfield{GPIO\_OUT\_W1TC}{32}{0}{{0x000000}}%
\reglabel{Reset}\regnewline%
\begin{regdesc}\begin{reglist}
\label{fielddesc:GPIOOUTW1TC}\item [GPIO\_OUT\_W1TC] 配置是否清零 GPIO0 \verb+~+ GPIO31 输出寄存器  \hyperref[regdesc:GPIOOUTREG]{GPIO\_OUT\_REG}。\\
0：不清零 \\
 1： \hyperref[regdesc:GPIOOUTREG]{GPIO\_OUT\_REG} 中相应位将被清零 \\
bit0 \verb+~+ bit31 分别对应 GPIO0 \verb+~+ GPIO31。
推荐使用此寄存器来清零  \hyperref[regdesc:GPIOOUTREG]{GPIO\_OUT\_REG}。\\
(WT)
\end{reglist}\end{regdesc}
\end{register}


\begin{register}{H}{GPIO\_OUT1\_REG}{0x0010}\label{regdesc:GPIOOUT1REG}
\regfield{(reserved)}{7}{25}{0000000}%
\regfield{GPIO\_OUT1\_DATA\_ORIG}{25}{0}{{0x00000}}%
\reglabel{Reset}\regnewline%
\begin{regdesc}\begin{reglist}
\label{fielddesc:GPIOOUT1DATAORIG}\item [GPIO\_OUT1\_DATA\_ORIG] 配置简单输出模式下 GPIO32 \verb+~+ GPIO54 的输出值。\\
0：低电平\\
1：高电平 \\
 bit0 \verb+~+ bit22 的值分别对应 GPIO32 \verb+~+ GPIO54 的输出值。\\
(R/W/SC/WTC)
\end{reglist}\end{regdesc}
\end{register}


\begin{register}{H}{GPIO\_OUT1\_W1TS\_REG}{0x0014}\label{regdesc:GPIOOUT1W1TSREG}
\regfield{(reserved)}{7}{25}{0000000}%
\regfield{GPIO\_OUT1\_W1TS}{25}{0}{{0x00000}}%
\reglabel{Reset}\regnewline%
\begin{regdesc}\begin{reglist}
\label{fielddesc:GPIOOUT1W1TS}\item [GPIO\_OUT1\_W1TS] 配置是否置位
 GPIO32 \verb+~+ GPIO54 的输出寄存器 \hyperref[regdesc:GPIOOUT1REG]{GPIO\_OUT1\_REG}。\\
0：不置位\\
1：\hyperref[regdesc:GPIOOUT1REG]{GPIO\_OUT1\_REG} 中相应位也将置 1 \\
bit0 \verb+~+ bit22 分别对应 GPIO32 \verb+~+ GPIO54。
推荐使用此寄存器来置位  \hyperref[regdesc:GPIOOUT1REG]{GPIO\_OUT1\_REG}。\\
(WT)
\end{reglist}\end{regdesc}
\end{register}


\begin{register}{H}{GPIO\_OUT1\_W1TC\_REG}{0x0018}\label{regdesc:GPIOOUT1W1TCREG}
\regfield{(reserved)}{7}{25}{0000000}%
\regfield{GPIO\_OUT1\_W1TC}{25}{0}{{0x00000}}%
\reglabel{Reset}\regnewline%
\begin{regdesc}\begin{reglist}
\label{fielddesc:GPIOOUT1W1TC}\item [GPIO\_OUT1\_W1TC] 配置是否清零 GPIO32 \verb+~+ GPIO54 输出寄存器  \hyperref[regdesc:GPIOOUT1REG]{GPIO\_OUT1\_REG}。\\
0：不清零 \\
 1： \hyperref[regdesc:GPIOOUT1REG]{GPIO\_OUT1\_REG} 中相应位将被清零 \\
bit0 \verb+~+ bit22 分别对应 GPIO32 \verb+~+ GPIO54。
推荐使用此寄存器来清零  \hyperref[regdesc:GPIOOUT1REG]{GPIO\_OUT\_REG}。\\
(WT)
\end{reglist}\end{regdesc}
\end{register}


\begin{register}{H}{GPIO\_ENABLE\_REG}{0x0020}\label{regdesc:GPIOENABLEREG}
\regfield{GPIO\_ENABLE\_DATA}{32}{0}{{0x000000}}%
\reglabel{Reset}\regnewline%
\begin{regdesc}\begin{reglist}
\label{fielddesc:GPIOENABLEDATA}\item [GPIO\_ENABLE\_DATA] 配置是否使能 GPIO0 \verb+~+ GPIO31 输出。\\
0：不使能\\
1：使能\\
bit0 \verb+~+ bit31 分别对应 GPIO0 \verb+~+ GPIO31。 \\
(R/W/WTC)
\end{reglist}\end{regdesc}
\end{register}


\begin{register}{H}{GPIO\_ENABLE\_W1TS\_REG}{0x0024}\label{regdesc:GPIOENABLEW1TSREG}
\regfield{GPIO\_ENABLE\_W1TS}{32}{0}{{0x000000}}%
\reglabel{Reset}\regnewline%
\begin{regdesc}\begin{reglist}
\label{fielddesc:GPIOENABLEW1TS}\item [GPIO\_ENABLE\_W1TS] 配置是否置位 GPIO0 \verb+~+ GPIO31 的输出使能寄存器 \hyperref[regdesc:GPIOENABLEREG]{GPIO\_ENABLE\_REG}。\\
0：不置位\\
1： \hyperref[regdesc:GPIOENABLEREG]{GPIO\_ENABLE\_REG} 中相应位也将置 1 \\
bit0 \verb+~+ bit31 分别对应 GPIO0 \verb+~+ GPIO31。推荐使用该寄存器来置位  \hyperref[regdesc:GPIOENABLEREG]{GPIO\_ENABLE\_REG}。\\
(WT)
\end{reglist}\end{regdesc}
\end{register}


\begin{register}{H}{GPIO\_ENABLE\_W1TC\_REG}{0x0028}\label{regdesc:GPIOENABLEW1TCREG}
\regfield{GPIO\_ENABLE\_W1TC}{32}{0}{{0x000000}}%
\reglabel{Reset}\regnewline%
\begin{regdesc}\begin{reglist}
\label{fielddesc:GPIOENABLEW1TC}\item [GPIO\_ENABLE\_W1TC] 配置是否清零 GPIO0 \verb+~+ GPIO31 的输出使能寄存器  \hyperref[regdesc:GPIOENABLEREG]{GPIO\_ENABLE\_REG}。\\
0：不清零 \\
 1： \hyperref[regdesc:GPIOENABLEREG]{GPIO\_ENABLE\_REG} 中相应位将被清零 \\
bit0 \verb+~+ bit31 分别对应 GPIO0 \verb+~+ GPIO31。推荐使用该寄存器清零 \hyperref[regdesc:GPIOENABLEREG]{GPIO\_ENABLE\_REG}。\\
(WT)
\end{reglist}\end{regdesc}
\end{register}


\begin{register}{H}{GPIO\_ENABLE1\_REG}{0x002C}\label{regdesc:GPIOENABLE1REG}
\regfield{(reserved)}{7}{25}{0000000}%
\regfield{GPIO\_ENABLE1\_DATA}{25}{0}{{0x00000}}%
\reglabel{Reset}\regnewline%
\begin{regdesc}\begin{reglist}
\label{fielddesc:GPIOENABLE1DATA}\item [GPIO\_ENABLE1\_DATA] 配置是否使能 GPIO32 \verb+~+ GPIO54 输出。\\
0：不使能\\
1：使能\\
bit0 \verb+~+ bit22 分别对应 GPIO32 \verb+~+ GPIO54。 \\
(R/W/WTC)
\end{reglist}\end{regdesc}
\end{register}


\begin{register}{H}{GPIO\_ENABLE1\_W1TS\_REG}{0x0030}\label{regdesc:GPIOENABLE1W1TSREG}
\regfield{(reserved)}{7}{25}{0000000}%
\regfield{GPIO\_ENABLE1\_W1TS}{25}{0}{{0x00000}}%
\reglabel{Reset}\regnewline%
\begin{regdesc}\begin{reglist}
\label{fielddesc:GPIOENABLE1W1TS}\item [GPIO\_ENABLE1\_W1TS] 配置是否置位 GPIO32 \verb+~+ GPIO54 的输出使能寄存器 \hyperref[regdesc:GPIOENABLE1REG]{GPIO\_ENABLE1\_REG}。\\
0：不置位\\
1： \hyperref[regdesc:GPIOENABLE1REG]{GPIO\_ENABLE1\_REG} 中相应位也将置 1 \\
bit0 \verb+~+ bit22 分别对应 GPIO32 \verb+~+ GPIO54。推荐使用该寄存器来置位  \hyperref[regdesc:GPIOENABLE1REG]{GPIO\_ENABLE1\_REG}。\\
(WT)
\end{reglist}\end{regdesc}
\end{register}


\begin{register}{H}{GPIO\_ENABLE1\_W1TC\_REG}{0x0034}\label{regdesc:GPIOENABLE1W1TCREG}
\regfield{(reserved)}{7}{25}{0000000}%
\regfield{GPIO\_ENABLE1\_W1TC}{25}{0}{{0x00000}}%
\reglabel{Reset}\regnewline%
\begin{regdesc}\begin{reglist}
\label{fielddesc:GPIOENABLE1W1TC}\item [GPIO\_ENABLE1\_W1TC] 配置是否清零 GPIO32 \verb+~+ GPIO54 的输出使能寄存器  \hyperref[regdesc:GPIOENABLE1REG]{GPIO\_ENABLE1\_REG}。\\
0：不清零 \\
 1：\hyperref[regdesc:GPIOENABLE1REG]{GPIO\_ENABLE1\_REG} 中相应位将被清零 \\
bit0 \verb+~+ bit22 分别对应 GPIO32 \verb+~+ GPIO54。推荐使用该寄存器清零  \hyperref[regdesc:GPIOENABLE1REG]{GPIO\_ENABLE1\_REG}。\\
(WT)
\end{reglist}\end{regdesc}
\end{register}


\begin{register}{H}{GPIO\_STRAP\_REG}{0x0038}\label{regdesc:GPIOSTRAPREG}
\regfield{(reserved)}{16}{16}{0000000000000000}%
\regfield{GPIO\_STRAPPING}{16}{0}{{0x00}}%
\reglabel{Reset}\regnewline%
\begin{regdesc}\begin{reglist}
\label{fielddesc:GPIOSTRAPPING}\item [GPIO\_STRAPPING] 表示 GPIO Strapping 管脚的值。\\
bit0：GPIO38\\
bit1：GPIO37\\
bit2：GPIO36\\
bit3：GPIO35\\
bit4：GPIO34\\
bit5 \verb+~+ bit15：无效\\
更多有关 Strapping 管脚控制的功能，见\textit{\href{https://www.espressif.com/sites/default/files/documentation/esp32-p4_datasheet_cn.pdf}{《\chipname{} 技术规格书》}} > 章节\textit{Strapping 管脚}。\\
(RO)
\end{reglist}\end{regdesc}
\end{register}

\begin{register}{H}{GPIO\_IN\_REG}{0x003C}\label{regdesc:GPIOINREG}
\regfield{GPIO\_IN\_DATA\_NEXT}{32}{0}{{0x000000}}%
\reglabel{Reset}\regnewline%
\begin{regdesc}\begin{reglist}
\label{fielddesc:GPIOINDATANEXT}\item [GPIO\_IN\_DATA\_NEXT] 表示 GPIO0 \verb+~+ GPIO31 的输入值。\\每一位均代表一个管脚的输入值：\\
0：低电平\\
1：高电平\\
bit0 \verb+~+ bit31 分别对应 GPIO0 \verb+~+ GPIO31。\\
(RO)
\end{reglist}\end{regdesc}
\end{register}


\begin{register}{H}{GPIO\_IN1\_REG}{0x0040}\label{regdesc:GPIOIN1REG}
\regfield{(reserved)}{7}{25}{0000000}%
\regfield{GPIO\_IN1\_DATA\_NEXT}{25}{0}{{0x00000}}%
\reglabel{Reset}\regnewline%
\begin{regdesc}\begin{reglist}
\label{fielddesc:GPIOIN1DATANEXT}\item [GPIO\_IN1\_DATA\_NEXT] 表示 GPIO32 \verb+~+ GPIO54 的输入值。每一位均代表一个管脚的输入值：\\
0：低电平\\
1：高电平\\
bit0 \verb+~+ bit22 分别对应 GPIO32 \verb+~+ GPIO54。 \\
(RO)
\end{reglist}\end{regdesc}
\end{register}


\begin{register}{H}{GPIO\_STATUS\_REG}{0x0044}\label{regdesc:GPIOSTATUSREG}
\regfield{GPIO\_STATUS\_INTERRUPT}{32}{0}{{0x000000}}%
\reglabel{Reset}\regnewline%
\begin{regdesc}\begin{reglist}
\label{fielddesc:GPIOSTATUSINTERRUPT}\item [GPIO\_STATUS\_INTERRUPT] GPIO0 \verb+~+ GPIO31 的中断状态寄存器，并且软件可配置该寄存器的值。\\
bit0 \verb+~+ bit31 分别对应 GPIO0 \verb+~+ GPIO31。\\
每一位均代表对应 GPIO 的状态：\\
0：表示对应 GPIO 未产生  \hyperref[fielddesc:GPIOPINNINTTYPE]{GPIO\_PIN\regindex{n}\_INT\_TYPE} 所选择的中断，或软件配置该位寄存器的值为 0。\\
1：表示对应 GPIO 产生了  \hyperref[fielddesc:GPIOPINNINTTYPE]{GPIO\_PIN\regindex{n}\_INT\_TYPE} 所选择的中断，或软件配置该位寄存器的值为 1。\\
(R/W/WTC)
\end{reglist}\end{regdesc}
\end{register}


\begin{register}{H}{GPIO\_STATUS\_W1TS\_REG}{0x0048}\label{regdesc:GPIOSTATUSW1TSREG}
\regfield{GPIO\_STATUS\_W1TS}{32}{0}{{0x000000}}%
\reglabel{Reset}\regnewline%
\begin{regdesc}\begin{reglist}
\label{fielddesc:GPIOSTATUSW1TS}\item [GPIO\_STATUS\_W1TS] 配置是否置位 GPIO0 \verb+~+ GPIO31 的中断状态寄存器 \hyperref[fielddesc:GPIOSTATUSINTERRUPT]{GPIO\_STATUS\_INTERRUPT}。\\
bit0 \verb+~+ bit31 分别对应 GPIO0 \verb+~+ GPIO31。\\
每一位置 1，则  \hyperref[fielddesc:GPIOSTATUSINTERRUPT]{GPIO\_STATUS\_INTERRUPT} 中的相应位也置 1。\\
推荐使用此寄存器来置位  \hyperref[fielddesc:GPIOSTATUSINTERRUPT]{GPIO\_STATUS\_INTERRUPT}。\\
(WT)
\end{reglist}\end{regdesc}
\end{register}


\begin{register}{H}{GPIO\_STATUS\_W1TC\_REG}{0x004C}\label{regdesc:GPIOSTATUSW1TCREG}
\regfield{GPIO\_STATUS\_W1TC}{32}{0}{{0x000000}}%
\reglabel{Reset}\regnewline%
\begin{regdesc}\begin{reglist}
\label{fielddesc:GPIOSTATUSW1TC}\item [GPIO\_STATUS\_W1TC] 配置是否清零 GPIO0 \verb+~+ GPIO31 的中断状态寄存器  \hyperref[fielddesc:GPIOSTATUSINTERRUPT]{GPIO\_STATUS\_INTERRUPT}。\\
bit0 \verb+~+ bit31 分别对应 GPIO0 \verb+~+ GPIO31。\\
每一位置 1，则  \hyperref[fielddesc:GPIOSTATUSINTERRUPT]{GPIO\_STATUS\_INTERRUPT} 中的相应位也清零。\\
推荐使用此寄存器来清零 \hyperref[fielddesc:GPIOSTATUSINTERRUPT]{GPIO\_STATUS\_INTERRUPT}。\\
(WT)
\end{reglist}\end{regdesc}
\end{register}


\begin{register}{H}{GPIO\_STATUS1\_REG}{0x0050}\label{regdesc:GPIOSTATUS1REG}
\regfield{(reserved)}{7}{25}{0000000}%
\regfield{GPIO\_STATUS1\_INTERRUPT}{25}{0}{{0x00000}}%
\reglabel{Reset}\regnewline%
\begin{regdesc}\begin{reglist}
\label{fielddesc:GPIOSTATUS1INTERRUPT}\item [GPIO\_STATUS1\_INTERRUPT] 表示 GPIO32 \verb+~+ GPIO54 的中断状态寄存器，并且软件可配置该寄存器的值。\\
bit0 \verb+~+ bit22 分别对应 GPIO32 \verb+~+ GPIO54。 \\
每一位均代表对应 GPIO 的状态：\\
0：表示对应 GPIO 未产生  \hyperref[fielddesc:GPIOPINNINTTYPE]{GPIO\_PIN\regindex{n}\_INT\_TYPE} 所选择的中断，或软件配置该位寄存器的值为 0。\\
1：表示对应 GPIO 产生了  \hyperref[fielddesc:GPIOPINNINTTYPE]{GPIO\_PIN\regindex{n}\_INT\_TYPE} 所选择的中断，或软件配置该位寄存器的值为 1。\\
(R/W/WTC)
\end{reglist}\end{regdesc}
\end{register}


\begin{register}{H}{GPIO\_STATUS1\_W1TS\_REG}{0x0054}\label{regdesc:GPIOSTATUS1W1TSREG}
\regfield{(reserved)}{7}{25}{0000000}%
\regfield{GPIO\_STATUS1\_W1TS}{25}{0}{{0x00000}}%
\reglabel{Reset}\regnewline%
\begin{regdesc}\begin{reglist}
\label{fielddesc:GPIOSTATUS1W1TS}\item [GPIO\_STATUS1\_W1TS] 配置是否置位 GPIO32 \verb+~+ GPIO54 的中断状态寄存器 \hyperref[fielddesc:GPIOSTATUS1INTERRUPT]{GPIO\_STATUS1\_INTERRUPT}。\\
bit0 \verb+~+ bit22 分别对应 GPIO32 \verb+~+ GPIO54。\\
每一位置 1，则  \hyperref[fielddesc:GPIOSTATUS1INTERRUPT]{GPIO\_STATUS1\_INTERRUPT} 中的相应位也置 1。\\
推荐使用此寄存器来置位  \hyperref[fielddesc:GPIOSTATUS1INTERRUPT]{GPIO\_STATUS1\_INTERRUPT}。\\
(WT)
\end{reglist}\end{regdesc}
\end{register}


\begin{register}{H}{GPIO\_STATUS1\_W1TC\_REG}{0x0058}\label{regdesc:GPIOSTATUS1W1TCREG}
\regfield{(reserved)}{7}{25}{0000000}%
\regfield{GPIO\_STATUS1\_W1TC}{25}{0}{{0x00000}}%
\reglabel{Reset}\regnewline%
\begin{regdesc}\begin{reglist}
\label{fielddesc:GPIOSTATUS1W1TC}\item [GPIO\_STATUS1\_W1TC] 配置是否清零 GPIO32 \verb+~+ GPIO54 的中断状态寄存器  \hyperref[fielddesc:GPIOSTATUS1INTERRUPT]{GPIO\_STATUS1\_INTERRUPT}。\\
bit0 \verb+~+ bit22 分别对应 GPIO32 \verb+~+ GPIO54。 \\
每一位置 1，则  \hyperref[fielddesc:GPIOSTATUS1INTERRUPT]{GPIO\_STATUS1\_INTERRUPT} 中的相应位将被清零。\\
推荐使用此寄存器来清零  \hyperref[fielddesc:GPIOSTATUS1INTERRUPT]{GPIO\_STATUS1\_INTERRUPT}。\\
(WT)
\end{reglist}\end{regdesc}
\end{register}


\begin{register}{H}{GPIO\_INTR\_0\_REG}{0x005C}\label{regdesc:GPIOINTR0REG}
\regfield{GPIO\_INT\_0}{32}{0}{{0x000000}}%
\reglabel{Reset}\regnewline%
\begin{regdesc}\begin{reglist}
\label{fielddesc:GPIOINT0}\item [GPIO\_INT\_0] 表示 GPIO0 \verb+~+ GPIO31 的 CPU 中断 0 的屏蔽中断状态位。 每位代表：\\
 0：表示 CPU 中断未使能，或对应 GPIO 未产生  \hyperref[fielddesc:GPIOPINNINTTYPE]{GPIO\_PIN\regindex{n}\_INT\_TYPE} 所配置的中断。\\
1：表示 CPU 中断使能后，对应 GPIO 产生了  \hyperref[fielddesc:GPIOPINNINTTYPE]{GPIO\_PIN\regindex{n}\_INT\_TYPE} 所配置的中断。\\
bit0 \verb+~+ bit31 分别对应 GPIO0 \verb+~+ GPIO31。如果  \hyperref[regdesc:GPIOPINNREG]{GPIO\_PIN\regindex{n}\_REG} 中 bit13 高电平有效，即使能 CPU 中断，则此寄存器所示的中断状态应与  \hyperref[regdesc:GPIOSTATUSREG]{GPIO\_STATUS\_REG} 中相应 bit 的中断状态一致。\\
(RO)
\end{reglist}\end{regdesc}
\end{register}


\begin{register}{H}{GPIO\_INTR1\_0\_REG}{0x0060}\label{regdesc:GPIOINTR10REG}
\regfield{(reserved)}{7}{25}{0000000}%
\regfield{GPIO\_INT1\_0}{25}{0}{{0x00000}}%
\reglabel{Reset}\regnewline%
\begin{regdesc}\begin{reglist}
\label{fielddesc:GPIOINT10}\item [GPIO\_INT1\_0] 表示 GPIO32 \verb+~+ GPIO54 的 CPU 中断 0 的屏蔽中断状态位。 每位代表：\\
 0：表示 CPU 中断未使能，或对应 GPIO 未产生  \hyperref[fielddesc:GPIOPINNINTTYPE]{GPIO\_PIN\regindex{n}\_INT\_TYPE} 所配置的中断。\\
1：表示 CPU 中断使能后，对应 GPIO 产生了  \hyperref[fielddesc:GPIOPINNINTTYPE]{GPIO\_PIN\regindex{n}\_INT\_TYPE} 所配置的中断。\\
bit0 \verb+~+ bit22 分别对应 GPIO32 \verb+~+ GPIO54。如果  \hyperref[regdesc:GPIOPINNREG]{GPIO\_PIN\regindex{n}\_REG} 中 bit13 高电平有效，即使能 CPU 中断，则此寄存器所示的中断状态应与  \hyperref[regdesc:GPIOSTATUS1REG]{GPIO\_STATUS1\_REG} 中相应 bit 的中断状态一致。\\
(RO)
\end{reglist}\end{regdesc}
\end{register}


\begin{register}{H}{GPIO\_INTR\_1\_REG}{0x0064}\label{regdesc:GPIOINTR1REG}
\regfield{GPIO\_INT\_1}{32}{0}{{0x000000}}%
\reglabel{Reset}\regnewline%
\begin{regdesc}\begin{reglist}
\label{fielddesc:GPIOINT1}\item [GPIO\_INT\_1] 表示 GPIO0 \verb+~+ GPIO31 的 CPU 中断 1 的屏蔽中断状态位。 每位代表：\\
 0：表示 CPU 中断未使能，或对应 GPIO 未产生  \hyperref[fielddesc:GPIOPINNINTTYPE]{GPIO\_PIN\regindex{n}\_INT\_TYPE} 所配置的中断。\\
1：表示 CPU 中断使能后，对应 GPIO 产生了  \hyperref[fielddesc:GPIOPINNINTTYPE]{GPIO\_PIN\regindex{n}\_INT\_TYPE} 所配置的中断。\\
bit0 \verb+~+ bit31 分别对应 GPIO0 \verb+~+ GPIO31。如果  \hyperref[regdesc:GPIOPINNREG]{GPIO\_PIN\regindex{n}\_REG} 中 bit14 高电平有效，即使能 CPU 中断，则此寄存器所示的中断状态应与  \hyperref[regdesc:GPIOSTATUSREG]{GPIO\_STATUS\_REG} 中相应 bit 的中断状态一致。\\
(RO)
\end{reglist}\end{regdesc}
\end{register}


\begin{register}{H}{GPIO\_INTR1\_1\_REG}{0x0068}\label{regdesc:GPIOINTR11REG}
\regfield{(reserved)}{7}{25}{0000000}%
\regfield{GPIO\_INT1\_1}{25}{0}{{0x00000}}%
\reglabel{Reset}\regnewline%
\begin{regdesc}\begin{reglist}
\label{fielddesc:GPIOINT11}\item [GPIO\_INT1\_1] 表示 GPIO32 \verb+~+ GPIO54 的 CPU 中断 1 的屏蔽中断状态位。 每位代表：\\
 0：表示 CPU 中断未使能，或对应 GPIO 未产生  \hyperref[fielddesc:GPIOPINNINTTYPE]{GPIO\_PIN\regindex{n}\_INT\_TYPE} 所配置的中断。\\
1：表示 CPU 中断使能后，对应 GPIO 产生了  \hyperref[fielddesc:GPIOPINNINTTYPE]{GPIO\_PIN\regindex{n}\_INT\_TYPE} 所配置的中断。\\
bit0 \verb+~+ bit22 分别对应 GPIO32 \verb+~+ GPIO54。如果  \hyperref[regdesc:GPIOPINNREG]{GPIO\_PIN\regindex{n}\_REG} 中 bit14 高电平有效，即使能 CPU 中断，则此寄存器所示的中断状态应与  \hyperref[regdesc:GPIOSTATUS1REG]{GPIO\_STATUS1\_REG} 中相应 bit 的中断状态一致。
(RO)
\end{reglist}\end{regdesc}
\end{register}


\begin{register}{H}{GPIO\_STATUS\_NEXT\_REG}{0x006C}\label{regdesc:GPIOSTATUSNEXTREG}
\regfield{GPIO\_STATUS\_INTERRUPT\_NEXT}{32}{0}{{0x000000}}%
\reglabel{Reset}\regnewline%
\begin{regdesc}\begin{reglist}
\label{fielddesc:GPIOSTATUSINTERRUPTNEXT}\item [GPIO\_STATUS\_INTERRUPT\_NEXT] 表示 GPIO0 \verb+~+ GPIO31 的中断源信号。\\
bit0 \verb+~+ bit31 分别对应 GPIO0 \verb+~+ 31。每位代表：\\
 0：表示 GPIO 未产生  \hyperref[fielddesc:GPIOPINNINTTYPE]{GPIO\_PIN\regindex{n}\_INT\_TYPE} 所配置的中断。\\
1：GPIO 产生了  \hyperref[fielddesc:GPIOPINNINTTYPE]{GPIO\_PIN\regindex{n}\_INT\_TYPE} 所配置的中断。\\
上述中断可以为上升沿中断、下降沿中断、电平敏感中断或任一沿中断。\\  (RO)
\end{reglist}\end{regdesc}
\end{register}


\begin{register}{H}{GPIO\_STATUS\_NEXT1\_REG}{0x0070}\label{regdesc:GPIOSTATUSNEXT1REG}
\regfield{(reserved)}{7}{25}{0000000}%
\regfield{GPIO\_STATUS\_INTERRUPT\_NEXT1}{25}{0}{{0x00000}}%
\reglabel{Reset}\regnewline%
\begin{regdesc}\begin{reglist}
\label{fielddesc:GPIOSTATUSINTERRUPTNEXT1}\item [GPIO\_STATUS\_INTERRUPT\_NEXT1] 表示 GPIO32 \verb+~+ GPIO54 的中断源信号。\\
bit0 \verb+~+ bit22 分别对应 GPIO32 \verb+~+ GPIO54。每位代表：\\
 0：表示 GPIO 未产生  \hyperref[fielddesc:GPIOPINNINTTYPE]{GPIO\_PIN\regindex{n}\_INT\_TYPE} 所配置的中断。\\
1：GPIO 产生了  \hyperref[fielddesc:GPIOPINNINTTYPE]{GPIO\_PIN\regindex{n}\_INT\_TYPE} 所配置的中断。\\
上述中断可以为上升沿中断、下降沿中断、电平敏感中断或任一沿中断。\\  (RO)
\end{reglist}\end{regdesc}
\end{register}


\begin{register}{H}{GPIO\_PIN\regindex{n}\_REG (\regindex{n}: 0 - 54)}
{0x0074+0x4*\regindex{n}}\label{regdesc:GPIOPINNREG}
\regfield{(reserved)}{14}{18}{00000000000000}%
\regfield{GPIO\_PIN\regindex{n}\_INT\_ENA}{5}{13}{{0x0}}%
\regfield{(reserved)}{2}{11}{{0x0}}%
\regfield{GPIO\_PIN\regindex{n}\_WAKEUP\_ENABLE}{1}{10}{{0}}%
\regfield{GPIO\_PIN\regindex{n}\_INT\_TYPE}{3}{7}{{0x0}}%
\regfield{(reserved)}{2}{5}{00}%
\regfield{GPIO\_PIN\regindex{n}\_SYNC1\_BYPASS}{2}{3}{{0x0}}%
\regfield{GPIO\_PIN\regindex{n}\_PAD\_DRIVER}{1}{2}{{0}}%
\regfield{GPIO\_PIN\regindex{n}\_SYNC2\_BYPASS}{2}{0}{{0x0}}%
\reglabel{Reset}\regnewline%
\begin{regdesc}\begin{reglist}
\label{fielddesc:GPIOPINNSYNC2BYPASS}\item [GPIO\_PIN\regindex{n}\_SYNC2\_BYPASS] 配置是否使能 GPIO 输入数据在第二级同步中 HP IO MUX 运行时钟上升沿同步或下降沿同步。 \\
0：不同步\\
1：在下降沿进行同步\\
2：在上升沿进行同步\\
3：在上升沿进行同步\\

(R/W)
\label{fielddesc:GPIOPINNPADDRIVER}\item [GPIO\_PIN\regindex{n}\_PAD\_DRIVER] 配置选择管脚的输出驱动模式。\\
0：正常输出\\
1：开漏输出\\(R/W)
\label{fielddesc:GPIOPINNSYNC1BYPASS}\item [GPIO\_PIN\regindex{n}\_SYNC1\_BYPASS] 配置是否使能 GPIO 输入数据在第一级同步中 HP IO MUX 运行时钟上升沿同步或下降沿同步。 \\
0：不同步\\
1：在下降沿进行同步\\
2：在上升沿进行同步\\
3：在上升沿进行同步\\
(R/W)
\label{fielddesc:GPIOPINNINTTYPE}\item [GPIO\_PIN\regindex{n}\_INT\_TYPE] 配置 GPIO 中断类型。 \\
0：关闭 GPIO 中断\\
1：上升沿触发\\
2：下降沿触发\\
3：任一沿触发\\
4：低电平触发\\
5：高电平触发\\ (R/W)

\label{fielddesc:GPIOPINNWAKEUPENABLE}\item [GPIO\_PIN\regindex{n}\_WAKEUP\_ENABLE] 配置是否使能 GPIO 唤醒功能。\\
0：不使能\\
1：使能\\
该功能只能将 CPU 从 Light-sleep 中唤醒。\\(R/W)
\item [见下页]
\end{reglist}\end{regdesc}
\end{register}

\addtocounter{Regfloat}{-1}
\begin{register}{H}{GPIO\_PIN\regindex{n}\_REG (\regindex{n}: 0 - 54)}{0x0074+0x4*\regindex{n}}
\begin{regdesc}\begin{reglist}
\item [接上页]
\label{fielddesc:GPIOPINNINTENA}\item [GPIO\_PIN\regindex{n}\_INT\_ENA] 配置是否使能 CPU 中断或 CPU 非屏蔽中断。\\
bit13：配置是否使能 CPU 中断 0：\\
0：不使能\\
1：使能\\
bit14：配置是否使能 CPU 中断 1：\\
0：不使能\\
1：使能\\
bit15：无效\\
bit16：配置是否使能 CPU 中断 2：\\
0：不使能\\
1：使能\\
bit17：配置是否使能 CPU 中断 3：\\
0：不使能\\
1：使能\\
(R/W)
\end{reglist}\end{regdesc}
\end{register}

\begin{register}{H}{GPIO\_FUNC\regindex{n}\_IN\_SEL\_CFG\_REG (\regindex{n}: 0 - 255)} {0x015C+4*\regindex{n}}\label{regdesc:GPIOFUNCNINSELCFGREG}
\regfield{(reserved)}{24}{8}{000000000000000000000000}%
\regfield{GPIO\_SIG\regindex{n}\_IN\_SEL}{1}{7}{{0}}%
\regfield{GPIO\_FUNC\regindex{n}\_IN\_INV\_SEL}{1}{6}{{0}}%
\regfield{GPIO\_FUNC\regindex{n}\_IN\_SEL}{6}{0}{{0x3f}}%
\reglabel{Reset}\regnewline%
\begin{regdesc}\begin{reglist}
\label{fielddesc:GPIOFUNCNINSEL}\item [GPIO\_FUNC\regindex{n}\_IN\_SEL] 配置从 55 个 GPIO 中选择一个管脚连接输入信号 \regindex{n}。\\
0：选择 GPIO0 \\
1：选择 GPIO1\\
......\\
53：选择 GPIO53 \\
54：选择 GPIO54\\
55 \verb+~+ 61：无效 \\
62：选择恒低电平输入\\
63：选择恒高电平输入\\

(R/W)
\label{fielddesc:GPIOFUNCNININVSEL}\item [GPIO\_FUNC\regindex{n}\_IN\_INV\_SEL] 配置是否反转输入值。\\
0：不反转\\
1：反转\\ (R/W)
\label{fielddesc:GPIOSIGNINSEL}\item [GPIO\_SIG\regindex{n}\_IN\_SEL] 配置是否通过 HP GPIO 交换矩阵连接信号。\\
0：旁路 HP GPIO 交换矩阵，即直接通过 HP IO MUX 连接信号与外设\\
1：通过 HP GPIO 交换矩阵连接信号与外设\\  (R/W)
\end{reglist}\end{regdesc}
\end{register}
\begin{register}{H}{GPIO\_FUNC\regindex{n}\_OUT\_SEL\_CFG\_REG (\regindex{n}: 0 - 54)}{0x0558+0x4*\regindex{n}}\label{regdesc:GPIOFUNCNOUTSELCFGREG}
\regfield{(reserved)}{20}{12}{00000000000000000000}%
\regfield{GPIO\_FUNC\regindex{n}\_OE\_INV\_SEL}{1}{11}{{0}}%
\regfield{GPIO\_FUNC\regindex{n}\_OE\_SEL}{1}{10}{{0}}%
\regfield{GPIO\_FUNC\regindex{n}\_OUT\_INV\_SEL}{1}{9}{{0}}%
\regfield{GPIO\_FUNC\regindex{n}\_OUT\_SEL}{9}{0}{{0x100}}%
\reglabel{Reset}\regnewline%
\begin{regdesc}\begin{reglist}
\label{fielddesc:GPIOFUNCNOUTSEL}\item [GPIO\_FUNC\regindex{n}\_OUT\_SEL] 配置从 256 个外设信号中选择一个信号  \regindex{Y} (0 <= \regindex{Y} < 256) 输出至 GPIO\regindex{n}。\\
0：选择信号 0\\
1：选择信号 1\\
......\\
254：选择信号 254\\
255：选择信号 255\\
256：\hyperref[regdesc:GPIOOUTREG]{GPIO\_OUT\_REG}/ \hyperref[regdesc:GPIOOUT1REG]{GPIO\_OUT1\_REG} 和  \hyperref[regdesc:GPIOENABLEREG]{GPIO\_ENABLE\_REG}/ \hyperref[regdesc:GPIOENABLE1REG]{GPIO\_ENABLE1\_REG} 中的 bit\regindex{n} 将用作输出值和输出使能信号。\\
257 \verb+~+ 511：无效\\

信号列表见表 \ref{tab:iomuxgpio-periph-signals-via-hp-gpio-matrix-cn}。

(R/W/SC)
\label{fielddesc:GPIOFUNCNOUTINVSEL}\item [GPIO\_FUNC\regindex{n}\_OUT\_INV\_SEL] 配置是否反转输出值。\\
0：不反转\\
1：反转\\ (R/W/SC)
\label{fielddesc:GPIOFUNCNOESEL}\item [GPIO\_FUNC\regindex{n}\_OE\_SEL] 配置选择输出使能信号。\\
0：使用外设的输出使能信号\\
 1：强制使用 \hyperref[regdesc:GPIOENABLEREG]{GPIO\_ENABLE\_REG} 的 bit\regindex{n} (\regindex{n}: 0 \verb+~+ 31) 和  \hyperref[regdesc:GPIOENABLE1REG]{GPIO\_ENABLE1\_REG} 的 bit\regindex{n} - 32 (\regindex{n}: 32 \verb+~+ 54) 用作输出使能信号。\\(R/W)
\label{fielddesc:GPIOFUNCNOEINVSEL}\item [GPIO\_FUNC\regindex{n}\_OE\_INV\_SEL] 配置是否反转输出使能信号。\\
0：不反转\\
1：反转\\ (R/W)
\end{reglist}\end{regdesc}
\end{register}


\begin{register}{H}{GPIO\_INTR\_2\_REG}{0x063C}\label{regdesc:GPIOINTR2REG}
\regfield{GPIO\_INT\_2}{32}{0}{{0x000000}}%
\reglabel{Reset}\regnewline%
\begin{regdesc}\begin{reglist}
\label{fielddesc:GPIOINT2}\item [GPIO\_INT\_2] 表示 GPIO0 \verb+~+ GPIO31 的 CPU 中断 2 的屏蔽中断状态位。 每位代表：\\
 0：表示 CPU 中断未使能，或对应 GPIO 未产生  \hyperref[fielddesc:GPIOPINNINTTYPE]{GPIO\_PIN\regindex{n}\_INT\_TYPE} 所配置的中断。\\
1：表示 CPU 中断使能后，对应 GPIO 产生了  \hyperref[fielddesc:GPIOPINNINTTYPE]{GPIO\_PIN\regindex{n}\_INT\_TYPE} 所配置的中断。\\
bit0 \verb+~+ bit31 分别对应 GPIO0 \verb+~+ GPIO31。如果  \hyperref[regdesc:GPIOPINNREG]{GPIO\_PIN\regindex{n}\_REG} 中 bit16 高电平有效，即使能 CPU 中断，则此寄存器所示的中断状态应与 \hyperref[regdesc:GPIOSTATUSREG]{GPIO\_STATUS\_REG} 中相应 bit 的中断状态一致。
(RO)
\end{reglist}\end{regdesc}
\end{register}


\begin{register}{H}{GPIO\_INTR1\_2\_REG}{0x0640}\label{regdesc:GPIOINTR12REG}
\regfield{(reserved)}{7}{25}{0000000}%
\regfield{GPIO\_INT1\_2}{25}{0}{{0x00000}}%
\reglabel{Reset}\regnewline%
\begin{regdesc}\begin{reglist}
\label{fielddesc:GPIOINT12}\item [GPIO\_INT1\_2] 表示 GPIO32 \verb+~+ GPIO54 的 CPU 中断 2 的屏蔽中断状态位。 每位代表：\\
 0：表示 CPU 中断未使能，或对应 GPIO 未产生  \hyperref[fielddesc:GPIOPINNINTTYPE]{GPIO\_PIN\regindex{n}\_INT\_TYPE} 所配置的中断。\\
1：表示 CPU 中断使能后，对应 GPIO 产生了  \hyperref[fielddesc:GPIOPINNINTTYPE]{GPIO\_PIN\regindex{n}\_INT\_TYPE} 所配置的中断。\\
bit0 \verb+~+ bit22 分别对应 GPIO32 \verb+~+ GPIO54。如果  \hyperref[regdesc:GPIOPINNREG]{GPIO\_PIN\regindex{n}\_REG} 中 bit16 高电平有效，即使能 CPU 中断，则此寄存器所示的中断状态应与  \hyperref[regdesc:GPIOSTATUS1REG]{GPIO\_STATUS1\_REG} 中相应 bit 的中断状态一致。\\
(RO)
\end{reglist}\end{regdesc}
\end{register}


\begin{register}{H}{GPIO\_INTR\_3\_REG}{0x0644}\label{regdesc:GPIOINTR3REG}
\regfield{GPIO\_INT\_3}{32}{0}{{0x000000}}%
\reglabel{Reset}\regnewline%
\begin{regdesc}\begin{reglist}
\label{fielddesc:GPIOINT3}\item [GPIO\_INT\_3] 表示 GPIO0 \verb+~+ GPIO31 的 CPU 中断 3 的屏蔽中断状态位。 每位代表：\\
 0：表示 CPU 中断未使能，或对应 GPIO 未产生  \hyperref[fielddesc:GPIOPINNINTTYPE]{GPIO\_PIN\regindex{n}\_INT\_TYPE} 所配置的中断。\\
1：表示 CPU 中断使能后，对应 GPIO 产生了  \hyperref[fielddesc:GPIOPINNINTTYPE]{GPIO\_PIN\regindex{n}\_INT\_TYPE} 所配置的中断。\\
bit0 \verb+~+ bit31 分别对应 GPIO0 \verb+~+ GPIO31。如果  \hyperref[regdesc:GPIOPINNREG]{GPIO\_PIN\regindex{n}\_REG} 中 bit17 高电平有效，即使能 CPU 中断，则此寄存器所示的中断状态应与  \hyperref[regdesc:GPIOSTATUSREG]{GPIO\_STATUS\_REG} 中相应 bit 的中断状态一致。\\
(RO)
\end{reglist}\end{regdesc}
\end{register}


\begin{register}{H}{GPIO\_INTR1\_3\_REG}{0x0648}\label{regdesc:GPIOINTR13REG}
\regfield{(reserved)}{7}{25}{0000000}%
\regfield{GPIO\_INT1\_3}{25}{0}{{0x00000}}%
\reglabel{Reset}\regnewline%
\begin{regdesc}\begin{reglist}
\label{fielddesc:GPIOINT13}\item [GPIO\_INT1\_3] 表示 GPIO32 \verb+~+ GPIO54 的 CPU 中断 3 的屏蔽中断状态位。 每位代表：\\
 0：表示 CPU 中断未使能，或对应 GPIO 未产生  \hyperref[fielddesc:GPIOPINNINTTYPE]{GPIO\_PIN\regindex{n}\_INT\_TYPE} 所配置的中断。\\
1：表示 CPU 中断使能后，对应 GPIO 产生了  \hyperref[fielddesc:GPIOPINNINTTYPE]{GPIO\_PIN\regindex{n}\_INT\_TYPE} 所配置的中断。\\
Bit0 \verb+~+ bit22 分别对应 GPIO32 \verb+~+ GPIO54。如果  \hyperref[regdesc:GPIOPINNREG]{GPIO\_PIN\regindex{n}\_REG} 中 bit17 高电平有效，即使能 CPU 中断，则此寄存器所示的中断状态应与  \hyperref[regdesc:GPIOSTATUS1REG]{GPIO\_STATUS1\_REG} 中相应 bit 的中断状态一致。\\
(RO)
\end{reglist}\end{regdesc}
\end{register}


\begin{register}{H}{GPIO\_CLOCK\_GATE\_REG}{0x064C}\label{regdesc:GPIOCLOCKGATEREG}
\regfield{(reserved)}{31}{1}{0000000000000000000000000000000}%
\regfield{GPIO\_CLK\_EN}{1}{0}{{1}}%
\reglabel{Reset}\regnewline%
\begin{regdesc}\begin{reglist}
\label{fielddesc:GPIOCLKEN}\item [GPIO\_CLK\_EN] 配置是否使能时钟门控。\\
0：不使能\\
1：此位置 1，则时钟自由运转。\\(R/W)
\end{reglist}\end{regdesc}
\end{register}

\begin{register}{H}{GPIO\_ZERO\_DET0\_FILTER\_CNT\_REG}{0x0710}\label{regdesc:GPIOZERODET0FILTERCNTREG}
\regfield{GPIO\_ZERO\_DET0\_FILTER\_CNT}{32}{0}{{0xffffffff}}%
\reglabel{Reset}\regnewline%
\begin{regdesc}\begin{reglist}
\label{fielddesc:GPIOZERODET0FILTERCNT}\item [GPIO\_ZERO\_DET0\_FILTER\_CNT] 配置模拟电压比较器 0 的过零检测的阈值。\hyperref[fielddesc:GPIOZERODET0FILTERCNT]{GPIO\_ZERO\_DET0\_FILTER\_CNT} 个 HP IO MUX 运行时钟周期内的中断将被 CPU 屏蔽。\\ (R/W)
\end{reglist}\end{regdesc}
\end{register}


\begin{register}{H}{GPIO\_ZERO\_DET1\_FILTER\_CNT\_REG}{0x0714}\label{regdesc:GPIOZERODET1FILTERCNTREG}
\regfield{GPIO\_ZERO\_DET1\_FILTER\_CNT}{32}{0}{{0xffffffff}}%
\reglabel{Reset}\regnewline%
\begin{regdesc}\begin{reglist}
\label{fielddesc:GPIOZERODET1FILTERCNT}\item [GPIO\_ZERO\_DET1\_FILTER\_CNT] 配置模拟电压比较器 1 的过零检测的阈值。\hyperref[fielddesc:GPIOZERODET1FILTERCNT]{GPIO\_ZERO\_DET1\_FILTER\_CNT} 个 HP IO MUX 运行时钟周期内的中断将被 CPU 屏蔽。\\ (R/W)
\end{reglist}\end{regdesc}
\end{register}


\begin{register}{H}{GPIO\_DATE\_REG}{0x07FC}\label{regdesc:GPIODATEREG}
\regfield{(reserved)}{4}{28}{0000}%
\regfield{GPIO\_DATE}{28}{0}{{0x230403}}%
\reglabel{Reset}\regnewline%
\begin{regdesc}\begin{reglist}
\label{fielddesc:GPIODATE}\item [GPIO\_DATE] 版本控制寄存器。\\
(R/W)
\end{reglist}\end{regdesc}
\end{register}

\begin{register}{H}{GPIO\_INT\_RAW\_REG}{0x0700}\label{regdesc:GPIOINTRAWREG}
\regfield{(reserved)}{26}{6}{00000000000000000000000000}%
\regfield{GPIO\_COMP1\_ALL\_INT\_RAW}{1}{5}{{0}}%
\regfield{GPIO\_COMP1\_POS\_INT\_RAW}{1}{4}{{0}}%
\regfield{GPIO\_COMP1\_NEG\_INT\_RAW}{1}{3}{{0}}%
\regfield{GPIO\_COMP0\_ALL\_INT\_RAW}{1}{2}{{0}}%
\regfield{GPIO\_COMP0\_POS\_INT\_RAW}{1}{1}{{0}}%
\regfield{GPIO\_COMP0\_NEG\_INT\_RAW}{1}{0}{{0}}%
\reglabel{Reset}\regnewline%
\begin{regdesc}\begin{reglist}
\label{fielddesc:GPIOCOMP0NEGINTRAW}\item [GPIO\_COMP0\_NEG\_INT\_RAW] \hyperref[int:GPIOCOMPnNEGINT]{GPIO\_COMP0\_NEG\_INT} 的原始中断状态。 (R/WTC/SS)
\label{fielddesc:GPIOCOMP0POSINTRAW}\item [GPIO\_COMP0\_POS\_INT\_RAW] \hyperref[int:GPIOCOMPnPOSINT]{GPIO\_COMP0\_POS\_INT} 的原始中断状态。 (R/WTC/SS)
\label{fielddesc:GPIOCOMP0ALLINTRAW}\item [GPIO\_COMP0\_ALL\_INT\_RAW] \hyperref[int:GPIOCOMPnALLINT]{GPIO\_COMP0\_ALL\_INT} 的原始中断状态。 (R/WTC/SS)
\label{fielddesc:GPIOCOMP1NEGINTRAW}\item [GPIO\_COMP1\_NEG\_INT\_RAW] \hyperref[int:GPIOCOMPnNEGINT]{GPIO\_COMP1\_NEG\_INT} 的原始中断状态。 (R/WTC/SS)
\label{fielddesc:GPIOCOMP1POSINTRAW}\item [GPIO\_COMP1\_POS\_INT\_RAW] \hyperref[int:GPIOCOMPnPOSINT]{GPIO\_COMP1\_POS\_INT} 的原始中断状态。 (R/WTC/SS)
\label{fielddesc:GPIOCOMP1ALLINTRAW}\item [GPIO\_COMP1\_ALL\_INT\_RAW] \hyperref[int:GPIOCOMPnALLINT]{GPIO\_COMP1\_ALL\_INT} 的原始中断状态。 (R/WTC/SS)
\end{reglist}\end{regdesc}
\end{register}
\begin{register}{H}{GPIO\_INT\_ST\_REG}{0x0704}\label{regdesc:GPIOINTSTREG}
\regfield{(reserved)}{26}{6}{00000000000000000000000000}%
\regfield{GPIO\_COMP1\_ALL\_INT\_ST}{1}{5}{{0}}%
\regfield{GPIO\_COMP1\_POS\_INT\_ST}{1}{4}{{0}}%
\regfield{GPIO\_COMP1\_NEG\_INT\_ST}{1}{3}{{0}}%
\regfield{GPIO\_COMP0\_ALL\_INT\_ST}{1}{2}{{0}}%
\regfield{GPIO\_COMP0\_POS\_INT\_ST}{1}{1}{{0}}%
\regfield{GPIO\_COMP0\_NEG\_INT\_ST}{1}{0}{{0}}%
\reglabel{Reset}\regnewline%
\begin{regdesc}\begin{reglist}
\label{fielddesc:GPIOCOMP0NEGINTST}\item [GPIO\_COMP0\_NEG\_INT\_ST] \hyperref[int:GPIOCOMPnNEGINT]{GPIO\_COMP0\_NEG\_INT} 的屏蔽中断状态。 (RO)
\label{fielddesc:GPIOCOMP0POSINTST}\item [GPIO\_COMP0\_POS\_INT\_ST] \hyperref[int:GPIOCOMPnPOSINT]{GPIO\_COMP0\_POS\_INT} 的屏蔽中断状态。 (RO)
\label{fielddesc:GPIOCOMP0ALLINTST}\item [GPIO\_COMP0\_ALL\_INT\_ST]  \hyperref[int:GPIOCOMPnALLINT]{GPIO\_COMP0\_ALL\_INT} 的屏蔽中断状态。 (RO)
\label{fielddesc:GPIOCOMP1NEGINTST}\item [GPIO\_COMP1\_NEG\_INT\_ST] \hyperref[int:GPIOCOMPnNEGINT]{GPIO\_COMP1\_NEG\_INT} 的屏蔽中断状态。 (RO)
\label{fielddesc:GPIOCOMP1POSINTST}\item [GPIO\_COMP1\_POS\_INT\_ST] \hyperref[int:GPIOCOMPnPOSINT]{GPIO\_COMP1\_POS\_INT} 的屏蔽中断状态。 (RO)
\label{fielddesc:GPIOCOMP1ALLINTST}\item [GPIO\_COMP1\_ALL\_INT\_ST] \hyperref[int:GPIOCOMPnALLINT]{GPIO\_COMP1\_ALL\_INT} 的屏蔽中断状态。 (RO)
\end{reglist}\end{regdesc}
\end{register}


\begin{register}{H}{GPIO\_INT\_ENA\_REG}{0x0708}\label{regdesc:GPIOINTENAREG}
\regfield{(reserved)}{26}{6}{00000000000000000000000000}%
\regfield{GPIO\_COMP1\_ALL\_INT\_ENA}{1}{5}{{1}}%
\regfield{GPIO\_COMP1\_POS\_INT\_ENA}{1}{4}{{1}}%
\regfield{GPIO\_COMP1\_NEG\_INT\_ENA}{1}{3}{{1}}%
\regfield{GPIO\_COMP0\_ALL\_INT\_ENA}{1}{2}{{1}}%
\regfield{GPIO\_COMP0\_POS\_INT\_ENA}{1}{1}{{1}}%
\regfield{GPIO\_COMP0\_NEG\_INT\_ENA}{1}{0}{{1}}%
\reglabel{Reset}\regnewline%
\begin{regdesc}\begin{reglist}
\label{fielddesc:GPIOCOMP0NEGINTENA}\item [GPIO\_COMP0\_NEG\_INT\_ENA] 写 1 使能 \hyperref[int:GPIOCOMPnNEGINT]{GPIO\_COMP0\_NEG\_INT} 中断。 (R/W)
\label{fielddesc:GPIOCOMP0POSINTENA}\item [GPIO\_COMP0\_POS\_INT\_ENA] 写 1 使能 \hyperref[int:GPIOCOMPnPOSINT]{GPIO\_COMP0\_POS\_INT} 中断。 (R/W)
\label{fielddesc:GPIOCOMP0ALLINTENA}\item [GPIO\_COMP0\_ALL\_INT\_ENA] 写 1 使能  \hyperref[int:GPIOCOMPnALLINT]{GPIO\_COMP0\_ALL\_INT} 中断。 (R/W)
\label{fielddesc:GPIOCOMP1NEGINTENA}\item [GPIO\_COMP1\_NEG\_INT\_ENA] 写 1 使能 \hyperref[int:GPIOCOMPnNEGINT]{GPIO\_COMP1\_NEG\_INT} 中断。 (R/W)
\label{fielddesc:GPIOCOMP1POSINTENA}\item [GPIO\_COMP1\_POS\_INT\_ENA] 写 1 使能 \hyperref[int:GPIOCOMPnPOSINT]{GPIO\_COMP1\_POS\_INT} 中断。 (R/W)
\label{fielddesc:GPIOCOMP1ALLINTENA}\item [GPIO\_COMP1\_ALL\_INT\_ENA] 写 1 使能 \hyperref[int:GPIOCOMPnALLINT]{GPIO\_COMP1\_ALL\_INT} 中断。 (R/W)
\end{reglist}\end{regdesc}
\end{register}


\begin{register}{H}{GPIO\_INT\_CLR\_REG}{0x070C}\label{regdesc:GPIOINTCLRREG}
\regfield{(reserved)}{26}{6}{00000000000000000000000000}%
\regfield{GPIO\_COMP1\_ALL\_INT\_CLR}{1}{5}{{0}}%
\regfield{GPIO\_COMP1\_POS\_INT\_CLR}{1}{4}{{0}}%
\regfield{GPIO\_COMP1\_NEG\_INT\_CLR}{1}{3}{{0}}%
\regfield{GPIO\_COMP0\_ALL\_INT\_CLR}{1}{2}{{0}}%
\regfield{GPIO\_COMP0\_POS\_INT\_CLR}{1}{1}{{0}}%
\regfield{GPIO\_COMP0\_NEG\_INT\_CLR}{1}{0}{{0}}%
\reglabel{Reset}\regnewline%
\begin{regdesc}\begin{reglist}
\label{fielddesc:GPIOCOMP0NEGINTCLR}\item [GPIO\_COMP0\_NEG\_INT\_CLR] 写 1 清除 \hyperref[int:GPIOCOMPnNEGINT]{GPIO\_COMP0\_NEG\_INT}。 (WT)
\label{fielddesc:GPIOCOMP0POSINTCLR}\item [GPIO\_COMP0\_POS\_INT\_CLR] 写 1 清除 \hyperref[int:GPIOCOMPnPOSINT]{GPIO\_COMP0\_POS\_INT}。 (WT)
\label{fielddesc:GPIOCOMP0ALLINTCLR}\item [GPIO\_COMP0\_ALL\_INT\_CLR] 写 1 清除  \hyperref[int:GPIOCOMPnALLINT]{GPIO\_COMP0\_ALL\_INT}。 (WT)
\label{fielddesc:GPIOCOMP1NEGINTCLR}\item [GPIO\_COMP1\_NEG\_INT\_CLR] 写 1 清除 \hyperref[int:GPIOCOMPnNEGINT]{GPIO\_COMP1\_NEG\_INT}。 (WT)
\label{fielddesc:GPIOCOMP1POSINTCLR}\item [GPIO\_COMP1\_POS\_INT\_CLR] 写 1 清除 \hyperref[int:GPIOCOMPnPOSINT]{GPIO\_COMP1\_POS\_INT}。 (WT)
\label{fielddesc:GPIOCOMP1ALLINTCLR}\item [GPIO\_COMP1\_ALL\_INT\_CLR] 写 1 清除 \hyperref[int:GPIOCOMPnALLINT]{GPIO\_COMP1\_ALL\_INT}。 (WT)
\end{reglist}\end{regdesc}
\end{register}


